\section{Rectangular Choi matrices and range bounds}

\begin{definition}[Matrix of a linear map in matrix-unit bases]
    \label{def:linear_map_matrix}
    \lean{Matrix.linearMapMatrix}
    \leanok
    Let $\mathcal L:\MN{d}\to\MN{e}$ be complex-linear. Its coefficient
    matrix $C(\mathcal L)\in M_{e^2\times d^2}(\C)$ in the matrix-unit bases is
    defined by
    \begin{align}
        C(\mathcal L)_{(a,b),(i,j)}
         & =(\mathcal L(E_{ij}))_{ab}.
        \label{eq:rectangular_choi_coefficient_matrix}
    \end{align}
\end{definition}

\begin{definition}[Unnormalized rectangular Choi matrix]
    \label{def:rectangular_choi}
    \lean{Matrix.rectangularChoi}
    \leanok
    \uses{def:linear_map_matrix}
    For a complex-linear map $\mathcal L:\MN{d}\to\MN{e}$, its unnormalized
    rectangular Choi matrix $J(\mathcal L)\in\MN{de}$ is the reshaping
    \begin{align}
        J(\mathcal L)_{(i,a),(j,b)}
         & =C(\mathcal L)_{(a,b),(i,j)}
        =(\mathcal L(E_{ij}))_{ab}.
        \label{eq:rectangular_choi_entries}
    \end{align}
    When $d=e$, this is $d$ times the normalized Choi matrix convention in
    \cite[Proposition~2.1]{Wolf2012Quantum}.
\end{definition}

\begin{theorem}[Frobenius norm under rectangular Choi reshaping]
    \label{thm:rectangular_choi_parseval}
    \lean{Matrix.rectangularChoi_frobenius_parseval}
    \leanok
    \uses{def:linear_map_matrix, def:rectangular_choi}
    For every complex-linear map $\mathcal L:\MN{d}\to\MN{e}$,
    \begin{align}
        \re\tr\!\left(J(\mathcal L)^\dagger J(\mathcal L)\right)
         & =
        \re\tr\!\left(C(\mathcal L)^\dagger C(\mathcal L)\right).
        \label{eq:rectangular_choi_parseval}
    \end{align}
\end{theorem}

\begin{proof}\leanok
    \uses{def:rectangular_choi}
    The two traces are the sums of the squared absolute values of their
    respective matrix entries. The bijection
    \begin{align}
        ((i,a),(j,b))\longmapsto((i,j),(a,b))
        \notag
    \end{align}
    identifies the two sums by~(\ref{eq:rectangular_choi_entries}).
\end{proof}

\begin{theorem}[Matrix-unit Parseval identity for a rectangular Choi matrix]
    \label{thm:rectangular_choi_parseval_matrix_units}
    \lean{Matrix.rectangularChoi_frobenius_parseval_matrix_units}
    \leanok
    \uses{def:rectangular_choi}
    For every complex-linear map $\mathcal L:\MN{d}\to\MN{e}$,
    \begin{align}
        \re\tr\!\left(J(\mathcal L)^\dagger J(\mathcal L)\right)
         & =
        \sum_{i,j=0}^{d-1}\sum_{a,b=0}^{e-1}
        \left|(\mathcal L(E_{ij}))_{ab}\right|^2.
        \label{eq:rectangular_choi_parseval_matrix_units}
    \end{align}
\end{theorem}

\begin{proof}\leanok
    \uses{def:linear_map_matrix, thm:rectangular_choi_parseval}
    Apply Theorem~\ref{thm:rectangular_choi_parseval} and substitute
    (\ref{eq:rectangular_choi_coefficient_matrix}) into the entrywise
    Frobenius-norm formula for $C(\mathcal L)$.
\end{proof}

\begin{definition}[Hilbert--Schmidt contraction]
    \label{def:rectangular_hilbert_schmidt_contraction}
    \lean{Matrix.IsHilbertSchmidtContraction}
    \leanok
    A complex-linear map $\mathcal L:\MN{d}\to\MN{e}$ is a
    \emph{Hilbert--Schmidt contraction} if, for every $X\in\MN{d}$,
    \begin{align}
        \re\tr\!\left((\mathcal L(X))^\dagger\mathcal L(X)\right)
         & \leq \re\tr(X^\dagger X).
        \label{eq:rectangular_hilbert_schmidt_contraction}
    \end{align}
\end{definition}

\begin{theorem}[Rectangular Choi norm bounded by the range dimension]
    \label{thm:rectangular_choi_rank_bound}
    \lean{Matrix.rectangularChoi_frobenius_sq_le_finrank_range}
    \leanok
    \uses{def:rectangular_choi, def:rectangular_hilbert_schmidt_contraction}
    If $\mathcal L:\MN{d}\to\MN{e}$ is a Hilbert--Schmidt contraction, then
    \begin{align}
        \left\|J(\mathcal L)\right\|_F^2
         & \leq\dim_{\C}\range\mathcal L.
        \label{eq:rectangular_choi_rank_bound}
    \end{align}
\end{theorem}

\begin{proof}\leanok
    \uses{def:linear_map_matrix, thm:rectangular_choi_parseval}
    Set $G=C(\mathcal L)^\dagger C(\mathcal L)$. The contraction inequality
    applied in the matrix-unit bases shows that every eigenvalue of the
    positive semidefinite matrix $G$ is at most one. Hence
    \begin{align}
        \tr G
         & \leq\#\{k:\lambda_k(G)\neq0\}
        =\rank C(\mathcal L)
        =\dim_{\C}\range\mathcal L.
        \notag
    \end{align}
    Theorem~\ref{thm:rectangular_choi_parseval} identifies the left-hand side
    with $\|J(\mathcal L)\|_F^2$.
\end{proof}
